\section{Algoritmos}

Estas notas estudian algoritmos.

\begin{definicion}{Algoritmo}{algoritmo}
Secuencia de pasos que transforman un conjunto de valores de \concept{entrada} (\textit{input}) en un conjunto de valores de \concept{salida} (\textit{output}).
\end{definicion}

Para muchos algoritmos se puede realizar una \concept{implementación}: escribir un programa, con algún lenguaje de programación y en alguna máquina, que ejecuta el algoritmo para entradas provistas. 

Sin embargo, el algoritmo como tal es un objeto matemático abstracto, independiente de su implementación, por lo que también se puede describir mediante:
\begin{itemize}
  \item Lenguaje natural (e.g., español, inglés)
  \item Lenguaje formal (e.g., cualquier lenguaje de programación)
  \item Seudocódigo (una mezcla de los dos anteriores)
  \item Como hardware (esto es difícil de imaginar en este punto, pero sí)
\end{itemize}

Para estudiar algoritmos, lo más conveniente es escribirlos en seudocódigo.

\setinputbase{fundamentos/algoritmo}
\inputlist{seudocodigo, refs}

